\subsection{Nilpotent invariance of effectivity for a specified local system}

\begin{lemma}[Compatible germs of a coherent module]
\label{lem:C-RD9-coherent}
Let $X$ be locally noetherian and let $\mathcal F$ be coherent.
A family $m_x\in\mathcal F_x$, compatible with every generization,
is the family of germs of a unique global section of $\mathcal F$.
\end{lemma}
\begin{proof}
Work on a noetherian affine open. Spread $m_x$ to a section on a
neighbourhood $V_x$ and delete the closures of all points in the finite
set $\operatorname{Ass}(\mathcal F)$ which do not generalize $x$.
Every associated point left in $V_x$ then generalizes $x$, so the
section has the prescribed value there. On overlaps, differences
vanish at every associated point of $\mathcal F$. Such a difference
is zero: a nonzero cyclic submodule of $\mathcal F$ would have an
associated prime in $\operatorname{Ass}(\mathcal F)$ at which it
remained nonzero. The local sections therefore glue. Uniqueness is
immediate from their germs.
\end{proof}

\begin{theorem}[C-RD9: nilpotents introduce no new obstruction in C.4]
\label{thm:C-RD9}
\status{Proved; independently audited with the original source.}
Let $S$ be locally noetherian, and fix a local system of abelian
schemes in the sense of C.4. This specified system is effective if
and only if its reduction on $S_{\mathrm{red}}$ is effective.
Consequently every such system is effective whenever
$S_{\mathrm{red}}$ locally admits a finite surjective geometrically
unibranch cover. No assumption that a prime is nilpotent on $S$ is
required.
\end{theorem}
\begin{proof}
Reduction proves necessity. For sufficiency it is enough to handle a
square-zero thickening $S_0\hookrightarrow S$ with coherent ideal
$I$, and a system whose reduction is induced by $A_0/S_0$ with the
specified local identifications. We may first work on an affine open
of $S$.

Consider the Zariski sheaf of isomorphism classes of actual abelian
scheme liftings of $A_0$, with their chosen identification on $S_0$.
It is a sheaf because an automorphism of an abelian scheme reducing
to the identity across a nilpotent ideal is the identity: the Hom
scheme is unramified. Compatible isomorphism classes therefore have
unique descent isomorphisms, and Zariski descent for schemes is
effective. This sheaf is locally nonempty. Indeed every prescribed
local $A(s)$, together with its reduction identification, spreads to
an open neighbourhood by finite presentation.

The square-zero deformation calculation for abelian schemes makes
this sheaf a torsor under the coherent module
\[
 \mathcal F
 =R^1\pi_*\mathcal Hom_{\mathcal O_{A_0}}
       (\Omega^1_{A_0/S_0},\pi^*I)
 =\operatorname{Lie}(A_0)\otimes
   \operatorname{Lie}(A_0^\vee)\otimes I,
 \qquad \pi:A_0\longrightarrow S_0.
\]
The lift-torsor calculation is explicitly recorded in the proof of
\cite[Lemma~24, p.590]{VasiuZink2010Purity}; existence over every affine square-zero thickening is also
proved over an arbitrary ring in
\cite[Chapter~IV, Lemma~2.8.1, pp.148--149]{Messing1972}. Neither
lifting statement imposes a characteristic restriction. Only
the difference-torsor calculation is needed here: local existence is
supplied by the given local system. The displayed tensor identity
uses invariant differentials and
$R^1\pi_*\mathcal O_{A_0}=\operatorname{Lie}(A_0^\vee)$, with
cohomology and base change.

Since $S_0$ is affine, $H^1(S_0,\mathcal F)=0$. The torsor has a
global section, represented by an actual abelian scheme $B/S$.
Compare $B$ at each local point with the prescribed $A(s)$. Their
difference is a well-defined element of $\mathcal F_s$ and these
elements are compatible with all generizations. By
Lemma~\ref{lem:C-RD9-coherent} they are the germs of a section
$m\in\Gamma(S_0,\mathcal F)$. Translate the lifting $B$ by this
section in the torsor of liftings. The resulting abelian scheme has
the required local isomorphisms. They are unique when required to
respect the reduction, and hence obey every prescribed generization
compatibility.

This proves the square-zero assertion on affine opens. On overlaps,
the resulting schemes have the prescribed compatible local
isomorphisms; Lemma~\ref{lem:C-RD8-sections}, applied to their Isom
scheme, globalizes those isomorphisms and verifies their cocycle.
Thus the assertion holds on all locally noetherian $S$.

Locally the nilradical is nilpotent. Apply the square-zero result
successively to its finite filtration by powers, and then glue the
constructions by the same uniqueness. This proves sufficiency for
$S_{\mathrm{red}}\hookrightarrow S$. The last assertion follows
from Theorem~\ref{thm:C-RD8}. Every lifting and gluing in this proof
is at a finite nilpotent stage and is an actual scheme; no assertion
that an arbitrary infinite formal deformation algebraizes is used.
\end{proof}

\begin{remark}[A genuinely nonreduced failure]
For the ring $R$ and local system of Theorem~\ref{thm:C-RD10}, base
change each local scheme and its generic identification along
$R\to R[\epsilon]/(\epsilon^2)$. This is a local system on a
nonreduced noetherian base. Any global realization would reduce to
the impossible realization over $R$, so it also fails to be effective.
\end{remark}
